\section{Complete fixed-exponent families and exact diagnostics}

\begin{theorem}
\label{thm:small}
For every $a\in\{2,3,4\}$, every $b\geq1$ and distinct primes $p,q,r$,
the graph $G_{p^a q^a r^b}$ is not Laplacian integral.
\end{theorem}
\begin{proof}
The cutoff values are $R(2)=0$, $R(3)=4/5$, and $R(4)=2$.
Theorem~\ref{thm:linear} settles all $b$ for $a=2,3$, and all $b>2$
for $a=4$. At $(a,b)=(4,1)$, $D=145$ is nonsquare.
At $(4,2)$ the symmetric cubic is
\[
f_{4,2}(x)=x^3-86x^2+2304x-18816.
\]
Its values at $0,1,2,3,4$ modulo $5$ are $4,3,1,4,3$, so it has
no root modulo $5$. A cubic over a field without a root is irreducible.
Since this cubic is monic, reduction modulo $5$ proves irreducibility over
$\mathbb Q$. Its real nonzero roots lift to noninteger graph eigenvalues.
\end{proof}

\begin{corollary}
\label{cor:five-six}
The same all-$b$ conclusion holds for $a=5,6$. Thus it holds for every
$a\in\{2,3,4,5,6\}$.
\end{corollary}
\begin{proof}
The second cutoff gives $S(5)=14/13$ and $S(6)=9/5$, leaving only
$b=1$ in each case. The discriminants are $221$ and $313$, respectively.
Since $14^2<221<15^2$ and $17^2<313<18^2$, both give irrational
antisymmetric eigenvalues. Every larger $b$ is covered by
Theorem~\ref{thm:second}.
\end{proof}

\begin{corollary}
\label{cor:seven-nine}
Every positive $b$ also gives a nonintegral spectrum for $a=7,8,9$.
Thus every repeated exponent $a\in\{2,3,4,5,6,7,8,9\}$ is settled.
\end{corollary}
\begin{proof}
The third cutoff gives $T(7)=6/5$, $T(8)=12/7$, $T(9)=70/31$.
Only four pairs remain under these bounds:
\[
(a,b)=(7,1),(8,1),(9,1),(9,2).
\]
Their discriminants are $421,545,685,1489$, respectively, between the
consecutive squares with lower square roots $20,23,26,38$.
Each therefore gives irrational antisymmetric eigenvalues, and
Theorem~\ref{thm:third} covers every larger $b$.
\end{proof}

For comparison, the complete divisor certificates before applying the
linear cutoff are given in Table~\ref{tab:divisors}. The constants factor
as $M=40=2^3\cdot5$, $189=3^3\cdot7$, and
$576=2^6\cdot3^2$. Testing every divisor $d\leq\sqrt M$ in
\eqref{eq:divisor-conditions} gives exactly the displayed admissible pairs.
The corresponding cubics in Table~\ref{tab:cubics} have no root modulo
the listed primes. These finite certificates supply a second proof using
the complete exceptional list. Three pairs are in the boundary family;
$(4,2)$ is the additional symmetric-block obstruction.

\begin{table}[ht]
\centering
\begin{tabular}{rrrrr}
\toprule
$a$ & $M$ & $(d,e)$ & $b$ & $\sqrt D$\\
\midrule
2 & 40 & $(1,40)$ & 3 & 13\\
3 & 189 & $(1,189)$ & 10 & 47\\
4 & 576 & $(6,96)$ & 2 & 18\\
4 & 576 & $(1,576)$ & 21 & 115\\
\bottomrule
\end{tabular}
\caption{Every admissible factor pair for $a=2,3,4$.}
\label{tab:divisors}
\end{table}

\begin{table}[ht]
\centering
\begin{tabular}{clc}
\toprule
$(a,b)$ & $f_{a,b}(x)$ & Witness prime\\
\midrule
$(2,3)$ & $x^3-47x^2+702x-3312$ & 5\\
$(3,10)$ & $x^3-187x^2+11220x-214200$ & 13\\
$(4,2)$ & $x^3-86x^2+2304x-18816$ & 5\\
$(4,21)$ & $x^3-485x^2+75492x-3721536$ & 5\\
\bottomrule
\end{tabular}
\caption{Root-free modular reductions for the exceptional cubics.}
\label{tab:cubics}
\end{table}

The separate diagnostic range $a=2,3,4$, $1\leq b\leq30$ contains
$90$ cubics, each with a noninteger root certified in an interval between
consecutive integers. The reducible cases in that range are $b=2,6$
for $a=2$, $b=3,12$ for $a=3$, and $b=4,20$ for $a=4$.
Each has a linear and an irreducible quadratic factor.
These observations do not require whole-cubic irreducibility.
Indeed, for every $a\geq1$,
\begin{align*}
f_{a,a}(x)&=(x-2a(a+1))
 \bigl(x^2-(5a^2+2a)x+6a^4+3a^3\bigr),\\
f_{a,a(a+1)}(x)&=(x-2a(a+1)^2)
 \bigl(x^2-(3a^3+4a^2+2a)x+2a^6+6a^5+7a^4+3a^3\bigr).
\end{align*}
Consequently no prime can make every diagnostic cubic irreducible: the
monic factors in the diagonal family retain positive degree modulo every
prime.

For the second family the quadratic discriminant is $a^2(a^4+4a+4)$.
For $a\geq3$ the radicand is between $a^4$ and $(a^2+1)^2$, since
$2a^2-4a-3=2(a-3)^2+8(a-3)+3>0$.
At $a=2$ it is $28$, between $5^2$ and $6^2$.
Thus the quadratic is irreducible for every $a\geq2$.

\begin{remark}
At $a=1,b=2$ the cubic is $(x-3)(x-6)(x-8)$, although the full graph
has irrational antisymmetric eigenvalues $(13\pm\sqrt{33})/2$.
A direct construction of its $10$ vertices gives
\[
\det(xI-L)=x(x-10)(x-9)^3(x-7)^2(x-4)(x^2-13x+34).
\]
Thus a noninteger cubic root is a sufficient condition for graph
nonintegrality; the converse cannot be used without the other blocks.
\end{remark}
